\documentclass[a4paper, 12pt]{article}
\usepackage{graphicx} % Required for inserting images
\usepackage{fullpage}
\usepackage{amsmath}
\usepackage[dvipsnames]{xcolor}
\usepackage{float}
\usepackage{geometry}
\usepackage{biblatex}
\geometry{margin=1in}
\usepackage{enumitem}
\usepackage{parskip}
\usepackage{pgfplots}
\pgfplotsset{compat=1.18}
\usepackage[normalem]{ulem}
\usepackage[version=4]{mhchem}
\usepackage{tikz}
\usepackage{microtype}
\usepackage{gensymb}
\usepackage{caption}
\usepackage{cancel}
\usepackage{nicefrac}

\newcommand{\lab}[1]{\: \text{#1}}
\newcommand{\degC}{$\degree$C \,}
\newcommand{\degF}{$\degree$F \,}
\newcommand{\R}{\left(0.0821 \: \frac{L \cdot atm}{mol \cdot \text{K}}\right)}
\newcommand{\cunits}{$\frac{J}{g \degree \text{C}}$}
\newcommand{\Hf}{$\Delta H_\text{f}$} % heat of formation
\newcommand{\mathHf}{\Delta H_\text{f}} %heat of formation in math mode

\usepackage{hyperref}
\hypersetup{
colorlinks = true,
linkcolor = teal,
citecolor=teal,         % Set color for citations
filecolor=teal,         % Set color for file links
urlcolor=teal           % Set color for URLs
}


\begin{document}

\begin{center}
\textbf{{\Large AP Chem summer homework}} \\
Emily Zhang \\
2025-07-31
\end{center}

\setcounter{section}{1}

\section{Math problems}

\begin{enumerate}[label=\alph*.]
    \item test
\end{enumerate}

\section{Significant figures}


\begin{enumerate}[label=\alph*.]
    \item 5
    \item 3
    \item 4
    \item 3
    \item 5
    \item 1
    \item 1
    \item 6
    \item 5
    \item 6
\end{enumerate}

\section{Chapter 1}

\begin{enumerate}[label=\alph*.]
\item draw this later
\item iii
\item 19\%
\item Dimensional analysis-
\begin{enumerate}[label=\roman*.]
\item hi
\item $10^3$; $10^2$; $10^{-3}$; $10^{-6}$; $10^{-9}$
\item $$\frac{330 \lab{minutes}}{1 \lab{day}} \cdot $$
\end{enumerate}
\item Element: oxygen; compound: water; homogeneous mixture: dish soap; heterogeneous mixture: oil in water; solution: saltwater
\end{enumerate}

\section{Chapter 2}

\begin{enumerate}[label=\alph*.]
\item The mass number is simply the sum of protons plus neutrons in an atom, while average atomic mass is the average of all the masses of all isotopes of an element.  
\item Mass can neither be created nor destroyed
\item
\begin{gather*}
2\ce{H2} + \ce{O2} \to 2\ce{H2O} \\
\end{gather*}
\item Molar mass of \ce{K2SO4}: 174.26 \lab{g/mol}
\begin{gather*}
\ce{K2}\text{---}2 \cdot 39.1 = 78.2 \lab{g/mol} \\
\ce{S}\text{---}32.06 \lab{g/mol} \\
\ce{O4}\text{---}4 \cdot 16 = 64 \lab{g/mol}
\end{gather*}
Assuming 1 mol, K makes up 78.2 of 174.26 grams
\[\frac{78.2}{174.26} \times 100\% = \boxed{44.88 \%}\]
\item Nucleus: center of the atom consisting of protons and neutrons; orbitals: areas where electrons are most likely to be found, organised by energy level; electrons: particles with negligible mass that are found in the orbitals surrounding the nucleus, adding negative charge; neutrons: a particle found in the nucleus with neutral charge holding protons together
\item Atoms with the same number of protons but different numbers of electrons (same charge, different mass)
\item C-14
\item Ionic bonding is the bonding of atoms through electrostatic attraction, a result of one atom taking electrons from another to fill both their valence shells; covalent bonding is the bonding of atoms through the sharing of electrons.
\item A chemical formula shows the number of each type of atom in a compound, and a structural formula creates a visual to show where the atoms are connected to each other in space.
\item ion: an atom with a positive or negative net charge due to an unequal number of protons and electrons; cation: an ion with positive charge; anion: an ion with negative charge.
\item Metals are characterised by high conductivity, lustre, and malleability. They lose electrons to nonmetals in ionic bonding. Nonmetals are found on the upper right side of the periodic table, and they take electrons during ionic bonding. They are dull, brittle, and poor conductors of heat and electricity. They also form covalent bonds with each other. Metalloids are elements that fall in between these two sets of characteristics. They are semiconductors and can form both ionic and covalent bonds.
\item Periods are the rows, and groups are the columns on the periodic table.
\item 
\end{enumerate}

\end{document}